\section{RQ1: Multi-Layer versus Single-Layer Discovery}
\label{sec:rq1}

\subsection{Method}

\paragraph{Corpus.} We built a ground-truth corpus of \data{rqone-cases}
cases: \data{rqone-cases-source} source cases (Go \data{rqone-cases-go},
Python \data{rqone-cases-python}, Java \data{rqone-cases-java}, C/C++
\data{rqone-cases-c}, JavaScript/TypeScript \data{rqone-cases-js}),
\data{rqone-cases-config} configuration and key-material cases,
\data{rqone-cases-deps} dependency cases, \data{rqone-cases-binary} binaries
and \data{rqone-cases-network} network endpoints served from containers. Each
case carries labels for every cryptographic asset it contains, by file, line
and algorithm family, \data{rqone-labels} in total. Beyond plain positives,
\data{rqone-kind-hard-positive} cases are hard positives (aliased imports,
wrappers, constants defined elsewhere, multi-line calls) and
\data{rqone-kind-hard-negative} are hard negatives (cryptography in comments
and strings, dead code, look-alike names). Labels record what the code does,
not what any tool detects; cases that \tool{} misses stay in the corpus.

\paragraph{Tools.} We compared \tool{} with every open-source scanner we could
run: CBOMkit, both its source analyzer (sonar-cryptography) and its
filesystem scanner (CBOMkit-theia)~\cite{cbomkit}, cryptoscan-pqc~\cite{cryptoscanpqc},
pqc-migration-audit~\cite{pqcaudit} and Semgrep with its cryptography
rules~\cite{semgrep}. Each runs once in a container pinned by digest and sees
a copy of the corpus without labels. We also planned to include a scanner
named Keylens but could not find one. Findings are mapped to algorithm families and
scored at two granularities: \emph{file} (file and family; fair to tools that
report no lines, and our headline) and \emph{line} (file, line and family). A
tool that does not cover a layer scores misses there; that is the point of
the question. An independent script recomputes every score from the
per-item outcomes.

\subsection{Results}

\begin{table*}[t]
\centering\small
% Generated by paper/gen.py from paper/data. Do not edit.
\begin{tabular}{llrrrrrr}
\toprule
Tool & Version & \multicolumn{3}{c}{File level} & \multicolumn{3}{c}{Line level} \\
\cmidrule(lr){3-5}\cmidrule(lr){6-8}
 &  & P & R & F1 & P & R & F1 \\
\midrule
\textbf{\tool{}} & \scriptsize this work & 0.997 & 0.975 & 0.986 & 0.994 & 0.972 & 0.983 \\
cryptoscan-pqc & \scriptsize 0.1.1 & 0.729 & 0.269 & 0.393 & 0.519 & 0.227 & 0.315 \\
CBOMkit (action) & \scriptsize v2.3.0 & 0.840 & 0.131 & 0.227 & 0.783 & 0.130 & 0.223 \\
pqc-migration-audit & \scriptsize 2b11fc59 & 0.704 & 0.119 & 0.203 & 0.406 & 0.108 & 0.170 \\
Semgrep & \scriptsize 1.184.0+rules@16b9e6d4 & 0.737 & 0.044 & 0.083 & 0.700 & 0.039 & 0.073 \\
CBOMkit-theia & \scriptsize latest@46a72ead & 0.833 & 0.016 & 0.031 & 0.000 & 0.000 & 0.000 \\
\bottomrule
\end{tabular}

\caption{Precision (P), recall (R) and F1 on all \data{rqone-cases} corpus
cases. File level matches file and algorithm family; line level also matches
the line.}
\label{tab:rq1-tools}
\end{table*}

\begin{table*}[t]
\centering\small\setlength{\tabcolsep}{4pt}
% Generated by paper/gen.py from paper/data. Do not edit.
\begin{tabular}{lrrrrrrr}
\toprule
Layer & Cases & \textbf{\tool{}} & cryptoscan-pqc & CBOMkit (action) & pqc-migration-audit & Semgrep & CBOMkit-theia \\
\midrule
source: Go & 38 & 0.98 & 0.33 & 0.00 & 0.22 & 0.08 & 0.00 \\
source: Python & 40 & 0.98 & 0.46 & 0.25 & 0.27 & 0.10 & 0.00 \\
source: Java & 26 & 0.93 & 0.45 & 1.00 & 0.21 & 0.17 & 0.00 \\
source: C/C++ & 24 & 0.97 & 0.25 & 0.00 & 0.19 & 0.00 & 0.00 \\
source: JS/TS & 25 & 0.94 & 0.41 & 0.00 & 0.00 & 0.00 & 0.00 \\
config & 24 & 1.00 & 0.05 & 0.00 & 0.00 & 0.00 & 0.09 \\
deps & 14 & 1.00 & 0.00 & 0.00 & 0.00 & 0.00 & 0.00 \\
binary & 6 & 1.00 & 0.30 & 0.00 & 0.00 & 0.00 & 0.00 \\
network & 5 & 0.95 & 0.00 & 0.00 & 0.00 & 0.00 & 0.00 \\
\midrule
all & 202 & 0.97 & 0.27 & 0.13 & 0.12 & 0.04 & 0.02 \\
\bottomrule
\end{tabular}

\caption{File-level recall by layer and source language. Only \tool{} probes
network endpoints.}
\label{tab:rq1-layers}
\end{table*}

Table~\ref{tab:rq1-tools} gives the overall scores. \tool{} finds
\data{rqone-pqscope-file-tp} of \data{rqone-file-labels} file-level assets
(recall \data{rqone-pqscope-file-r}) with \data{rqone-pqscope-file-fp} false
positive (precision \data{rqone-pqscope-file-p}). The highest-recall
alternative, \data{rqone-best-competitor}, reaches recall
\data{rqone-best-competitor-r} at precision
\data{rqone-best-competitor-p}, so \tool{} finds \data{rqone-recall-ratio} as
many assets.

Table~\ref{tab:rq1-layers} shows where the difference comes from. The other
tools read source code, and \data{rqone-nonsource-share} of the labeled
assets (\data{rqone-nonsource-labels}) are in configuration, dependencies,
binaries and endpoints. The best competitor recall in those layers is
\data{rqone-competitors-config-maxr} for configuration,
\data{rqone-competitors-deps-maxr} for dependencies and
\data{rqone-competitors-binary-maxr} for binaries; \tool{}'s are
\data{rqone-pqscope-file-config-r}, \data{rqone-pqscope-file-deps-r} and
\data{rqone-pqscope-file-binary-r}. Within source code the comparison is
closer for the languages a competitor targets: on Java,
\data{rqone-best-java-competitor} reaches \data{rqone-best-java-competitor-r},
above \tool{}'s \data{rqone-pqscope-file-java-r}. Across languages \tool{}'s
source recall ranges from \data{rqone-pqscope-lang-min-r}
(\data{rqone-pqscope-lang-min}) to \data{rqone-pqscope-lang-max-r}
(\data{rqone-pqscope-lang-max}). All \data{rqone-pqscope-misses} of its
file-level misses are in hard-positive cases. Of these,
\data{rqone-pqscope-misses-source} are in source code where the algorithm name
is computed at run time or reached through indirection (string concatenation,
template literals, computed properties, parameters, function pointers, a
dynamic import, a type-only reference), which needs data-flow analysis that
\tool{} does not do. The last is the cipher suite of an endpoint that
answered with a HelloRetryRequest.

\begin{table}[t]
\centering\small
% Generated by paper/gen.py from paper/data. Do not edit.
\begin{tabular}{lrrrr}
\toprule
Configuration & P & R & F1 & $\Delta$R (pp) \\
\midrule
all layers & 0.997 & 0.975 & 0.986 & -- \\
without source & 1.000 & 0.366 & 0.535 & 60.9 \\
without config & 0.996 & 0.800 & 0.887 & 17.5 \\
without deps & 0.997 & 0.909 & 0.951 & 6.6 \\
without binary & 0.997 & 0.912 & 0.953 & 6.2 \\
without network & 0.997 & 0.912 & 0.953 & 6.2 \\
\bottomrule
\end{tabular}

\caption{Layer ablation for \tool{} at file level. $\Delta$R is the recall
lost, in percentage points, when one layer is turned off.}
\label{tab:rq1-ablation}
\end{table}

\paragraph{Which layers contribute.} Table~\ref{tab:rq1-ablation} turns off
one layer at a time. Source code matters most: without it recall falls by
\data{rqone-abl-without-source-drop} points. Configuration and key material
come next (\data{rqone-abl-without-config-drop} points), then dependencies
(\data{rqone-abl-without-deps-drop}), binaries
(\data{rqone-abl-without-binary-drop}) and the network
(\data{rqone-abl-without-network-drop}). Every layer finds assets no other
layer finds; none is redundant on this corpus.

\paragraph{Caveats.} The corpus and \tool{}'s rules share an author, so
\tool{}'s scores measure consistency with the labeling convention more than
independent accuracy; the competitors' scores depend on that convention too,
through the family mapping. The per-item outcomes are released so every
match can be checked. The corpus layer mix is a design choice, and a
source-heavy corpus would narrow the gap; the per-layer table lets readers
reweight.
